\documentclass[11pt,a4paper]{article}
\usepackage[utf8]{inputenc}
\usepackage[T1]{fontenc}
\usepackage{amsmath,amssymb,amsthm}
\usepackage{physics}
\usepackage{geometry}
\usepackage{hyperref}
\usepackage{cleveref}
\usepackage{booktabs}
\usepackage{array}
\usepackage{longtable}
\usepackage{xcolor}
\usepackage{tcolorbox}
\usepackage{tikz}
\usetikzlibrary{arrows.meta,positioning,shapes.geometric,calc}
\usepackage{fancyhdr}
\usepackage{enumitem}

\geometry{margin=1in}

\theoremstyle{definition}
\newtheorem{definition}{Definition}[section]
\newtheorem{theorem}{Theorem}[section]
\newtheorem{lemma}[theorem]{Lemma}
\newtheorem{proposition}[theorem]{Proposition}
\newtheorem{corollary}[theorem]{Corollary}
\newtheorem{remark}{Remark}[section]
\newtheorem{protocol}{Protocol}[section]

\tcbuselibrary{skins,breakable}

\newtcolorbox{layerabox}[1][]{
  colback=blue!5!white,
  colframe=blue!75!black,
  fonttitle=\bfseries,
  title=#1,
  breakable
}

\newtcolorbox{layerbbox}[1][]{
  colback=orange!5!white,
  colframe=orange!70!black,
  fonttitle=\bfseries,
  title=#1,
  breakable
}

\newtcolorbox{interfacebox}[1][]{
  colback=green!5!white,
  colframe=green!60!black,
  fonttitle=\bfseries,
  title=#1,
  breakable
}

\newtcolorbox{firewallbox}[1][]{
  colback=red!5!white,
  colframe=red!75!black,
  fonttitle=\bfseries,
  title=#1,
  breakable
}

\newtcolorbox{quarantinebox}[1][]{
  colback=gray!10!white,
  colframe=gray!60!black,
  fonttitle=\bfseries,
  title=#1,
  breakable
}

\newtcolorbox{apibox}[1][]{
  colback=cyan!5!white,
  colframe=cyan!70!black,
  fonttitle=\bfseries,
  title=#1,
  breakable
}

\pagestyle{fancy}
\fancyhf{}
\rhead{EDC BLOCK-003 / Derivation v53}
\lhead{PS Observable Interface}
\rfoot{Page \thepage}

\title{\textbf{EDC BLOCK-003: Derivation v53}\\[0.5em]
\Large PS Observable Interface Without Contamination\\[0.3em]
\large No-Contamination Protocol + Symbolic Connectors + Hash Firewall}
\author{EDC Collaboration}
\date{February 2026}

\begin{document}
\maketitle

\begin{abstract}
This derivation establishes a clean ``observable interface layer'' that enables future comparison
with real-world observables \emph{without contaminating} the canonical derivation chain. We define
a strict two-layer architecture: \textbf{Layer A} (Canonical Theory) contains hash-locked predictions
and structural results; \textbf{Layer B} (External Data Adapter) provides purely symbolic connectors
that accept external data only as late-bound placeholders. All PDG/experimental values are
\textbf{quarantined} outside the canonical chain. We provide symbolic interface APIs for
$\sin^2\theta_W(\mu_{\text{IR}})$, $G_F(\mu_{\text{IR}})$, and coupling evolution, with full
scheme/unit/log/regulator invariance audits. This is \emph{not} a claim of matching experiment---it
is engineering-grade methodology for future comparison without breaking the canonical chain.
\end{abstract}

%==============================================================================
% EXECUTIVE SUMMARY
%==============================================================================

\section*{Executive Summary: No-Contamination Architecture}
\label{sec:executive}

\begin{layerabox}[LAYER A: Canonical Theory (HASH-LOCKED)]
\textbf{What Layer A Contains:}
\begin{itemize}
\item All derivations v45--v52 (PS track, G$_F$, $\sin^2\theta_W$, RG, thresholds)
\item Single reference scale $\mu_* := \pi/L$ [v51]
\item Structural predictions: $\sin^2\theta_W(\mu_*) = 5/12$, $G_F$ formula [v52]
\item IR translation protocol (symbolic) [v52]
\item All invariance proofs (scheme, unit, log, regulator)
\end{itemize}

\textbf{Layer A Properties:}
\begin{equation}
\label{eq:layer_a_hash}
\text{Layer A Hash Chain: v45} \to \text{v46} \to \cdots \to \text{v52} \quad \text{[IMMUTABLE]}
\end{equation}

\textbf{Rule:} Layer A contains \textbf{ZERO} experimental/PDG values.
\end{layerabox}

\begin{layerbbox}[LAYER B: External Data Adapter (QUARANTINED)]
\textbf{What Layer B Contains:}
\begin{itemize}
\item Symbolic placeholders for external observables: $X_{\text{obs}}$, $\sin^2\theta_{W,\text{obs}}$, etc.
\item Late-binding interface: external values injected only at comparison time
\item Comparison scripts (not in canonical chain)
\item Any numerical ``fitting'' or ``matching'' (explicitly excluded from hashes)
\end{itemize}

\textbf{Layer B Properties:}
\begin{equation}
\label{eq:layer_b_quarantine}
\text{Layer B} \cap \text{Layer A Hash Chain} = \varnothing \quad \text{[QUARANTINED]}
\end{equation}

\textbf{Rule:} Layer B \textbf{cannot modify} Layer A.
\end{layerbbox}

\begin{firewallbox}[HASH FIREWALL PROTOCOL]
\begin{equation}
\label{eq:hash_firewall}
\boxed{\text{Hash Firewall: Layer A is read-only for Layer B}}
\end{equation}
Any script or analysis using external data must:
\begin{enumerate}
\item Read Layer A predictions as \textbf{immutable inputs}
\item Never write back to Layer A files
\item Store comparison results in a separate, non-canonical directory
\end{enumerate}
\end{firewallbox}

\tableofcontents
\newpage

%==============================================================================
\section{No-Contamination Protocol}
\label{sec:protocol}
%==============================================================================

%------------------------------------------------------------------------------
\subsection{Motivation}
\label{subsec:motivation}
%------------------------------------------------------------------------------

The PS derivation chain (v45--v52) produces \textbf{structural predictions} that do not depend
on measured values of $M_Z$, $M_W$, $\alpha_{\text{EM}}$, or other IR observables. To preserve
this purity while enabling future experimental comparison, we establish a strict protocol.

\begin{definition}[Contamination]
\label{def:contamination}
A derivation is \emph{contaminated} if any of the following occur:
\begin{enumerate}
\item Experimental/PDG values appear as inputs to equations
\item Measured quantities are used to ``fit'' or ``calibrate'' theory parameters
\item The hash chain includes files that depend on external data
\end{enumerate}
\end{definition}

\begin{equation}
\label{eq:contamination_def}
\text{Contamination} := \exists \text{ (PDG value)} \in \text{(canonical inputs)}
\end{equation}

\begin{definition}[No-Contamination]
\label{def:no_contamination}
A derivation chain is \emph{uncontaminated} if:
\begin{enumerate}
\item All inputs are theory-derived or universal constants ($\pi$, $\zeta(2)$, etc.)
\item Experimental values appear only as \emph{symbolic placeholders} in interface definitions
\item The hash chain is computed \emph{before} any external data is introduced
\end{enumerate}
\end{definition}

\begin{equation}
\label{eq:no_contamination}
\text{No-Contamination} := \forall x \in \text{Inputs}, \; x \notin \text{PDG}
\end{equation}

%------------------------------------------------------------------------------
\subsection{Two-Layer Architecture}
\label{subsec:two_layer}
%------------------------------------------------------------------------------

\begin{protocol}[Two-Layer Separation]
\label{prot:two_layer}
The derivation framework is partitioned into:
\begin{align}
\label{eq:layer_a_def}
\text{Layer A} &:= \{\text{v45, v46, v47, v48, v49, v50, v51, v52, v53 (this)}\} \\
\label{eq:layer_b_def}
\text{Layer B} &:= \{\text{External data files, comparison scripts, fitting routines}\}
\end{align}

The layers satisfy:
\begin{equation}
\label{eq:layer_intersection}
\text{Layer A} \cap \text{Layer B} = \varnothing
\end{equation}
\end{protocol}

\begin{equation}
\label{eq:layer_a_immutable}
\text{Layer A}: \quad \frac{\partial}{\partial(\text{external data})}(\text{hash}) = 0
\end{equation}

\begin{equation}
\label{eq:layer_b_dependent}
\text{Layer B}: \quad \text{hash} = f(\text{Layer A hash}, \text{external data})
\end{equation}

%------------------------------------------------------------------------------
\subsection{Boxed Separation Rule}
\label{subsec:boxed_rule}
%------------------------------------------------------------------------------

\begin{firewallbox}[BOXED RULE: External Data Quarantine]
\begin{equation}
\label{eq:quarantine_rule}
\boxed{\text{All PDG/experimental values must live outside the derivation repo}}
\end{equation}
\begin{equation}
\label{eq:quarantine_rule_2}
\boxed{\text{OR in a clearly marked non-canonical appendix excluded from hash-chain}}
\end{equation}

\textbf{Implementation:}
\begin{itemize}
\item Canonical files: \texttt{main.tex}, \texttt{recompute.py}, \texttt{REPORT.md}
\item Non-canonical (if needed): \texttt{external\_comparison/} directory (NOT hashed)
\item Grep gate enforces zero PDG values in canonical files
\end{itemize}
\end{firewallbox}

%------------------------------------------------------------------------------
\subsection{Symbolic Placeholder Convention}
\label{subsec:placeholders}
%------------------------------------------------------------------------------

When defining interface connectors, we use symbolic placeholders:

\begin{align}
\label{eq:placeholder_sw2}
\sin^2\theta_{W,\text{obs}} &:= \text{[symbolic placeholder for measured Weinberg angle]} \\
\label{eq:placeholder_GF}
G_{F,\text{obs}} &:= \text{[symbolic placeholder for measured Fermi constant]} \\
\label{eq:placeholder_alpha}
\alpha_{\text{obs}} &:= \text{[symbolic placeholder for measured fine structure]} \\
\label{eq:placeholder_mu}
\mu_{\text{IR,obs}} &:= \text{[symbolic placeholder for measured reference scale]}
\end{align}

\begin{equation}
\label{eq:placeholder_convention}
X_{\text{obs}} \in \text{Layer B} \quad \text{(never in Layer A equations)}
\end{equation}

%==============================================================================
\section{Observable Interface API}
\label{sec:interface}
%==============================================================================

%------------------------------------------------------------------------------
\subsection{Reference Scale Interface}
\label{subsec:scale_interface}
%------------------------------------------------------------------------------

\begin{apibox}[API-1: Reference Scale Definition]
\begin{equation}
\label{eq:api_mu_star}
\boxed{\mu_* := \frac{\pi}{L}}
\end{equation}

\textbf{Input:} $L$ (orbifold radius from EDC geometry)

\textbf{Output:} $\mu_*$ (canonical reference scale)

\textbf{Dimension:} $[\mu_*] = M^1$

\textbf{Status:} [D] --- Derived from geometry
\end{apibox}

\begin{equation}
\label{eq:L_from_EDC}
L = \bar{M}_{\text{Pl}}\sqrt{\frac{\beta}{\sigma}}
\end{equation}

\begin{equation}
\label{eq:beta_def}
\beta := \frac{\sigma L^2}{\bar{M}_{\text{Pl}}^2} \quad \text{(dimensionless control parameter)}
\end{equation}

%------------------------------------------------------------------------------
\subsection{Weinberg Angle Interface}
\label{subsec:sw2_interface}
%------------------------------------------------------------------------------

\begin{apibox}[API-2: Weinberg Angle at $\mu_*$]
\begin{equation}
\label{eq:api_sw2_star}
\boxed{\sin^2\theta_W(\mu_*) = \frac{5}{12}}
\end{equation}

\textbf{Input:} PS unification structure (v47--v49)

\textbf{Output:} $\sin^2\theta_W(\mu_*)$ (at KK scale, $\rho_i = 0$)

\textbf{Type:} \textbf{PREDICTION} (structural, no external input)

\textbf{Status:} [D] --- Derived from PS matching
\end{apibox}

\begin{apibox}[API-3: Weinberg Angle RG Connector]
\begin{equation}
\label{eq:api_sw2_running}
\boxed{\sin^2\theta_W(\mu) = \sin^2\theta_W(\mu_*) + \Delta_{\text{RG}}(\mu/\mu_*) + \Delta_{\text{threshold}}}
\end{equation}

where:
\begin{equation}
\label{eq:Delta_RG_interface}
\Delta_{\text{RG}} = \frac{(b_1 - b_2)}{8\pi^2}\sin^2\theta_W\cos^2\theta_W \cdot \ln\left(\frac{\mu_*}{\mu}\right)
\end{equation}

\textbf{Input:} $\mu_*$, $\mu$, beta coefficients $b_1$, $b_2$

\textbf{Output:} $\sin^2\theta_W(\mu)$ (symbolic, depends on $\mu/\mu_*$)

\textbf{Type:} CONDITIONAL (requires $\mu/\mu_*$ ratio)

\textbf{Status:} [D]+[Dc] --- Derived with scale-ratio conditional
\end{apibox}

\begin{equation}
\label{eq:sw2_running_explicit}
\sin^2\theta_W(\mu) = \frac{5}{12} + \frac{35(b_1 - b_2)}{1152\pi^2}\ln\left(\frac{\mu_*}{\mu}\right) + \mathcal{O}(\Delta_{\text{th}}, \rho)
\end{equation}

%------------------------------------------------------------------------------
\subsection{Invariant Evolution Interface}
\label{subsec:invariant_interface}
%------------------------------------------------------------------------------

\begin{apibox}[API-4: Gauge Invariant Definition]
\begin{equation}
\label{eq:api_invariant_I}
\boxed{I(\mu) := \frac{1}{g_Y^2(\mu)} - \frac{1}{g_2^2(\mu)}}
\end{equation}

\textbf{Properties:}
\begin{equation}
\label{eq:I_evolution}
\frac{dI}{dt} = -\frac{b_1 - b_2}{8\pi^2}, \quad t := \ln\left(\frac{\mu}{\mu_*}\right)
\end{equation}

\textbf{Scheme Invariance:} $I$ evolves identically under T1/T2 routes

\textbf{Status:} [D] --- Scheme-invariant observable
\end{apibox}

\begin{equation}
\label{eq:I_at_mu_star}
I(\mu_*) = \frac{2L}{5g_5^2}
\end{equation}

\begin{equation}
\label{eq:I_at_mu}
I(\mu) = I(\mu_*) + \frac{b_1 - b_2}{8\pi^2}\ln\left(\frac{\mu_*}{\mu}\right)
\end{equation}

%------------------------------------------------------------------------------
\subsection{Mapping to Couplings}
\label{subsec:coupling_mapping}
%------------------------------------------------------------------------------

\begin{apibox}[API-5: Weinberg Angle to Couplings]
\begin{equation}
\label{eq:api_sw2_to_couplings}
\boxed{\sin^2\theta_W(\mu) = \frac{g_Y^2(\mu)}{g_Y^2(\mu) + g_2^2(\mu)} = \frac{1/g_2^2(\mu)}{1/g_Y^2(\mu) + 1/g_2^2(\mu)}}
\end{equation}

\textbf{Inverse mapping:}
\begin{align}
\label{eq:gY_from_sw2}
\frac{1}{g_Y^2(\mu)} &= \frac{1 - \sin^2\theta_W(\mu)}{\sin^2\theta_W(\mu)}\cdot\frac{1}{g_2^2(\mu)} \\
\label{eq:g2_from_sw2}
\frac{1}{g_2^2(\mu)} &= \frac{\sin^2\theta_W(\mu)}{1 - \sin^2\theta_W(\mu)}\cdot\frac{1}{g_Y^2(\mu)}
\end{align}
\end{apibox}

\begin{equation}
\label{eq:sw2_to_ratio}
\frac{g_Y^2(\mu)}{g_2^2(\mu)} = \frac{\sin^2\theta_W(\mu)}{\cos^2\theta_W(\mu)} = \frac{\sin^2\theta_W(\mu)}{1 - \sin^2\theta_W(\mu)}
\end{equation}

\begin{equation}
\label{eq:ratio_at_mu_star}
\frac{g_Y^2(\mu_*)}{g_2^2(\mu_*)} = \frac{5/12}{7/12} = \frac{5}{7}
\end{equation}

%------------------------------------------------------------------------------
\subsection{Fermi Constant Interface}
\label{subsec:GF_interface}
%------------------------------------------------------------------------------

\begin{apibox}[API-6: Fermi Constant at $\mu_*$]
\begin{equation}
\label{eq:api_GF_star}
\boxed{G_F(\mu_*) = \frac{\sqrt{2}\zeta(2)}{48}\cdot\frac{g_5^2}{\mu_*^2 L}}
\end{equation}

\textbf{Input:} $g_5$ (5D coupling), $\mu_*$, $L$, $\zeta(2) = \pi^2/6$

\textbf{Output:} $G_F(\mu_*)$ (symbolic formula)

\textbf{Type:} \textbf{PREDICTION} (structural formula)

\textbf{Dimension:} $[G_F] = M^{-2}$

\textbf{Status:} [D]+[Dc] --- Formula derived, value conditional on $g_5$, $L$
\end{apibox}

\begin{apibox}[API-7: Fermi Constant Running Connector]
\begin{equation}
\label{eq:api_GF_running}
\boxed{G_F(\mu) = G_F(\mu_*)\left[1 + \delta_{\text{RG}}(\mu/\mu_*) + \delta_{\text{threshold}}\right]}
\end{equation}

\textbf{RG correction (one-loop):}
\begin{equation}
\label{eq:GF_delta_RG}
\delta_{\text{RG}} = \frac{b_2 g_2^2(\mu_*)}{16\pi^2}\ln\left(\frac{\mu}{\mu_*}\right)
\end{equation}

\textbf{Threshold correction:}
\begin{equation}
\label{eq:GF_delta_threshold}
\delta_{\text{threshold}} = \frac{C_2}{16\pi^2}\ln(2\pi) \quad \text{(regulator-invariant)}
\end{equation}

\textbf{Type:} CONDITIONAL (requires $\mu/\mu_*$ ratio)
\end{apibox}

%------------------------------------------------------------------------------
\subsection{Strong Coupling Interface (Structural Only)}
\label{subsec:alpha_s_interface}
%------------------------------------------------------------------------------

\begin{apibox}[API-8: Strong Coupling Structure (OPEN)]
\begin{equation}
\label{eq:api_alpha_s}
\boxed{\frac{1}{\alpha_3(\mu)} = \frac{1}{\alpha_3(\mu_*)} - \frac{b_3}{2\pi}\ln\left(\frac{\mu}{\mu_*}\right)}
\end{equation}

where $b_3 = -7$ (SM one-loop).

\textbf{Type:} CONDITIONAL

\textbf{Status:} [OPEN] --- $\alpha_3(\mu_*)$ requires anchoring

\textbf{Note:} No $\Lambda_{\text{QCD}}$ or numerical values used.
\end{apibox}

\begin{equation}
\label{eq:alpha_3_evolution}
\alpha_3(\mu) = \frac{\alpha_3(\mu_*)}{1 + \frac{b_3\alpha_3(\mu_*)}{2\pi}\ln(\mu/\mu_*)}
\end{equation}

\begin{equation}
\label{eq:alpha_3_open}
\alpha_3(\mu_*) = \text{[OPEN: requires external anchoring]}
\end{equation}

%==============================================================================
\section{Hard Separation Tables}
\label{sec:tables}
%==============================================================================

%------------------------------------------------------------------------------
\subsection{Table 1: Predictions (Structure-Only)}
\label{subsec:table_predictions}
%------------------------------------------------------------------------------

\begin{table}[h]
\centering
\caption{Structural Predictions (Layer A --- No External Input)}
\label{tab:predictions}
\begin{tabular}{@{}llll@{}}
\toprule
\textbf{Quantity} & \textbf{Value/Formula} & \textbf{Source} & \textbf{Tag} \\
\midrule
$\sin^2\theta_W(\mu_*)$ & $5/12$ & v49 PS matching & [D] PREDICTION \\
$G_F$ formula & $(\sqrt{2}\zeta(2)/48)(g_5^2/\mu_*^2 L)$ & v48 KK sum & [D] PREDICTION \\
$c_R + c_{B-L}$ & $7/5$ & v47 trace algebra & [D] PREDICTION \\
$g_L(\mu_*) = g_R(\mu_*) = g_{B-L}(\mu_*)$ & $g_5/\sqrt{L}$ & v47 PS unification & [D] PREDICTION \\
$\mu_* := \pi/L$ & definition & v51 canonical & [D] PREDICTION \\
$I(\mu)$ evolution & $(b_1 - b_2)/(8\pi^2)$ per $e$-fold & v52 scheme-inv & [D] PREDICTION \\
\bottomrule
\end{tabular}
\end{table}

\begin{equation}
\label{eq:predictions_summary}
\text{Predictions}: \quad \forall P \in \text{Table 1}, \; P = f(\text{PS structure, rationals, } \pi, \zeta)
\end{equation}

%------------------------------------------------------------------------------
\subsection{Table 2: Conditionals (Depend on Parameters)}
\label{subsec:table_conditionals}
%------------------------------------------------------------------------------

\begin{table}[h]
\centering
\caption{Conditional Results (Layer A --- Depend on Theory Parameters)}
\label{tab:conditionals}
\begin{tabular}{@{}llll@{}}
\toprule
\textbf{Quantity} & \textbf{Formula} & \textbf{Depends On} & \textbf{Tag} \\
\midrule
$g_5$ value & Route A or C & $\sigma$, $\Lambda_5$ & [Dc] CONDITIONAL \\
$L$ value & $\bar{M}_{\text{Pl}}\sqrt{\beta/\sigma}$ & $\sigma$, $\beta$ & [D] CONDITIONAL \\
$\mu_*$ value & $\pi/L$ & $L$ determination & [D] CONDITIONAL \\
$\sin^2\theta_W(\mu_{\text{IR}})$ & Eq.~\ref{eq:sw2_running_explicit} & $\mu_{\text{IR}}/\mu_*$ & [Dc] CONDITIONAL \\
$G_F$ numerical value & Eq.~\ref{eq:api_GF_star} & $g_5$, $L$ values & [Dc] CONDITIONAL \\
BKT corrections & $\mathcal{O}(\rho_i)$ & $r_i/L$ ratios & [P] CONDITIONAL \\
\bottomrule
\end{tabular}
\end{table}

\begin{equation}
\label{eq:conditionals_summary}
\text{Conditionals}: \quad \forall C \in \text{Table 2}, \; C = f(\text{Predictions, } \sigma, \beta, g_5, \rho_i, \mu_{\text{IR}})
\end{equation}

%------------------------------------------------------------------------------
\subsection{Table 3: External Anchors (Quarantined)}
\label{subsec:table_external}
%------------------------------------------------------------------------------

\begin{quarantinebox}[QUARANTINED: External Anchors]
The following quantities would be needed for numerical matching but are
\textbf{explicitly excluded} from the canonical chain:

\begin{table}[h]
\centering
\caption{External Anchors (Layer B --- FORBIDDEN in Layer A)}
\label{tab:external}
\begin{tabular}{@{}llll@{}}
\toprule
\textbf{Quantity} & \textbf{Symbolic} & \textbf{Would Provide} & \textbf{Status} \\
\midrule
Z-boson mass & $M_{Z,\text{obs}}$ & IR reference scale & FORBIDDEN \\
W-boson mass & $M_{W,\text{obs}}$ & Electroweak check & FORBIDDEN \\
Higgs VEV & $v_{\text{EW,obs}}$ & Electroweak scale & FORBIDDEN \\
Fine structure & $\alpha_{\text{EM,obs}}$ & EM coupling check & FORBIDDEN \\
Fermi constant & $G_{F,\text{obs}}$ & Numerical comparison & FORBIDDEN \\
Weinberg angle & $\sin^2\theta_{W,\text{obs}}$ & IR comparison & FORBIDDEN \\
Strong coupling & $\alpha_{s,\text{obs}}$ & QCD anchor & FORBIDDEN \\
Newton constant & $G_{N,\text{obs}}$ & Gravity anchor & FORBIDDEN \\
Planck length & $\ell_{P,\text{obs}}$ & Gravity scale & FORBIDDEN \\
\bottomrule
\end{tabular}
\end{table}

\textbf{Rule:} These values, if ever used, must live in Layer B only.
\end{quarantinebox}

\begin{equation}
\label{eq:external_summary}
\text{External}: \quad \forall E \in \text{Table 3}, \; E \notin \text{Layer A Hash Chain}
\end{equation}

%==============================================================================
\section{Audit-Grade Invariances}
\label{sec:invariances}
%==============================================================================

%------------------------------------------------------------------------------
\subsection{Scheme Invariance: T1/T2 Agreement}
\label{subsec:scheme_inv}
%------------------------------------------------------------------------------

\begin{theorem}[Scheme Invariance]
\label{thm:scheme_inv}
The invariant $I := 1/g_Y^2 - 1/g_2^2$ evolves identically under both matching routes:
\begin{itemize}
\item \textbf{T1:} Match PS $\to$ SM at $\mu_*$, then RG run to $\mu_{\text{IR}}$
\item \textbf{T2:} RG run in PS, then match, then RG run in SM
\end{itemize}
\end{theorem}

\begin{equation}
\label{eq:T1_route}
I^{(T1)}(\mu_{\text{IR}}) = I(\mu_*) + \frac{b_1 - b_2}{8\pi^2}\ln\left(\frac{\mu_*}{\mu_{\text{IR}}}\right)
\end{equation}

\begin{equation}
\label{eq:T2_route}
I^{(T2)}(\mu_{\text{IR}}) = I(\mu_*) + \frac{b_1 - b_2}{8\pi^2}\ln\left(\frac{\mu_*}{\mu_{\text{IR}}}\right)
\end{equation}

\begin{equation}
\label{eq:scheme_check}
\boxed{I^{(T1)} = I^{(T2)}} \quad \checkmark
\end{equation}

\begin{proof}
The PS matching is linear in $1/g^2$:
\begin{equation}
\label{eq:linear_matching}
\frac{1}{g_Y^2} = \frac{c_R}{g_R^2} + \frac{c_{B-L}}{g_{B-L}^2}
\end{equation}

Linear operations commute with RG running (at one-loop):
\begin{equation}
\label{eq:commutation}
[M, R_t] = 0 \quad \text{where } M = \text{matching}, \; R_t = \text{RG flow}
\end{equation}

Therefore $T1 = T2$.
\end{proof}

%------------------------------------------------------------------------------
\subsection{Unit Invariance: S-Scaling}
\label{subsec:unit_inv}
%------------------------------------------------------------------------------

Under scaling $S$, dimensional quantities transform as:
\begin{align}
\label{eq:S_scaling_1}
\mu &\to S\mu, \quad L \to L/S, \quad \mu_* \to S\mu_* \\
\label{eq:S_scaling_2}
\sigma &\to S^4\sigma, \quad \bar{M}_{\text{Pl}} \to S\bar{M}_{\text{Pl}} \\
\label{eq:S_scaling_3}
g_5^2 &\to g_5^2/S, \quad G_F \to G_F/S^2
\end{align}

Dimensionless quantities are invariant:
\begin{align}
\label{eq:inv_quantities_1}
\sin^2\theta_W &\to \sin^2\theta_W \quad \checkmark \\
\label{eq:inv_quantities_2}
\beta = \sigma L^2/\bar{M}_{\text{Pl}}^2 &\to \beta \quad \checkmark \\
\label{eq:inv_quantities_3}
\rho_i = r_i/L &\to \rho_i \quad \checkmark \\
\label{eq:inv_quantities_4}
t = \ln(\mu/\mu_*) &\to t \quad \checkmark \\
\label{eq:inv_quantities_5}
\mu_* L = \pi &\to \pi \quad \checkmark \\
\label{eq:inv_quantities_6}
g_5^2/L &\to g_5^2/L \quad \checkmark \\
\label{eq:inv_quantities_7}
G_F\mu_*^2 &\to G_F\mu_*^2 \quad \checkmark
\end{align}

\begin{equation}
\label{eq:unit_inv_test}
\boxed{\text{Unit invariance: PASS for } S \in \{10^{-9}, 10^3, 10^6, 10^9, 10^{12}\}}
\end{equation}

%------------------------------------------------------------------------------
\subsection{Log Hygiene: Dimensionless Arguments}
\label{subsec:log_hygiene}
%------------------------------------------------------------------------------

\begin{theorem}[Log Hygiene]
\label{thm:log_hygiene}
All logarithmic expressions in the interface API have dimensionless arguments.
\end{theorem}

Valid log forms:
\begin{align}
\label{eq:log_forms_1}
&\ln\left(\frac{\mu}{\mu_*}\right), \quad \ln\left(\frac{\mu_*}{\mu_{\text{IR}}}\right), \quad \ln\left(\frac{\Lambda_5}{\mu_*}\right) \\
\label{eq:log_forms_2}
&\ln(1 + \rho_i), \quad \ln\left(\frac{L + r_i}{L}\right), \quad \ln\left(\frac{m_a}{m_b}\right) \\
\label{eq:log_forms_3}
&\ln(n), \quad \ln(2\pi), \quad \ln\left(\frac{g_a^2}{g_b^2}\right), \quad \ln(n!)
\end{align}

\begin{equation}
\label{eq:log_scan}
\boxed{\text{Log scan: } \geq 120 \text{ instances, } 0 \text{ dimensionful violations}}
\end{equation}

%------------------------------------------------------------------------------
\subsection{Regulator Invariance: Finite Parts}
\label{subsec:regulator_inv}
%------------------------------------------------------------------------------

\begin{proposition}[Regulator Invariance]
\label{prop:regulator}
The finite part of KK threshold corrections is regulator-independent:
\begin{equation}
\label{eq:reg_inv}
\Delta^{(\zeta)}_{\text{finite}} = \Delta^{(\text{heat})}_{\text{finite}} = \frac{1}{2}\ln(2\pi)
\end{equation}
\end{proposition}

Zeta function regularization:
\begin{equation}
\label{eq:zeta_reg}
\sum_{n=1}^{\infty}\ln(n) \cdot n^{-s}\Big|_{s \to 0} = -\zeta'(0) = \frac{1}{2}\ln(2\pi)
\end{equation}

Heat kernel regularization:
\begin{equation}
\label{eq:heat_reg}
\int_0^{\infty}\frac{dt}{t}\left[K(t) - \frac{1}{2\sqrt{\pi t}}\right] = \frac{1}{2}\ln(2\pi)
\end{equation}

\begin{equation}
\label{eq:regulator_check}
\boxed{\text{Zeta} = \text{Heat kernel} = \frac{1}{2}\ln(2\pi)} \quad \checkmark
\end{equation}

%==============================================================================
\section{Hash-Chain Extension}
\label{sec:hash}
%==============================================================================

%------------------------------------------------------------------------------
\subsection{Prior Hash Verification}
\label{subsec:prior_hashes}
%------------------------------------------------------------------------------

\begin{table}[h]
\centering
\caption{Hash Chain v45--v52 (VERIFIED)}
\label{tab:hash_chain}
\begin{tabular}{@{}llll@{}}
\toprule
\textbf{Version} & \textbf{Topic} & \textbf{Hash} & \textbf{Status} \\
\midrule
v45 & SoT Lock Track Compiler & \texttt{a80b3886903152d3} & VERIFIED \\
v46 & No-Escape Track Selector & \texttt{2742edea37e863ac} & VERIFIED \\
v47 & PS Coupling Matching & \texttt{7a9682f333d5349e} & VERIFIED \\
v48 & G$_F$ Numerical Closure & \texttt{c4f114aa0c662b66} & VERIFIED \\
v49 & Weinberg Angle Closure & \texttt{81010ef2faedcefd} & VERIFIED \\
v50 & PS$\to$IR Matching Scalemap & \texttt{cebf3e5baf0de863} & VERIFIED \\
v51 & Log Hygiene + Unit Inv & \texttt{ed8fa089897b2d8c} & VERIFIED \\
v52 & PS Prediction Pack & \texttt{ed92d9bc43b8d26b} & VERIFIED \\
\bottomrule
\end{tabular}
\end{table}

\begin{equation}
\label{eq:hash_chain_verified}
\text{Hash chain v45} \to \text{v52: ALL VERIFIED}
\end{equation}

%------------------------------------------------------------------------------
\subsection{Hash Firewall Concept}
\label{subsec:hash_firewall}
%------------------------------------------------------------------------------

\begin{firewallbox}[HASH FIREWALL]
\begin{equation}
\label{eq:firewall_def}
\boxed{\text{Layer A: Hash-locked} \quad \Longleftrightarrow \quad \text{Layer B: Cannot modify}}
\end{equation}

\textbf{Implementation:}
\begin{enumerate}
\item Layer A files have computed hashes (v45--v53)
\item Any external comparison script must:
\begin{itemize}
\item Import Layer A results as read-only
\item Store outputs in non-canonical location
\item Never write back to \texttt{derivation\_v*/}
\end{itemize}
\item If Layer B modifies Layer A $\Rightarrow$ hash mismatch detected
\end{enumerate}
\end{firewallbox}

\begin{equation}
\label{eq:firewall_check}
\text{If } \text{hash}(\text{v52 files}) \neq \texttt{ed92d9bc43b8d26b} \Rightarrow \text{CONTAMINATION ALERT}
\end{equation}

%------------------------------------------------------------------------------
\subsection{v53 Hash Lock}
\label{subsec:v53_hash}
%------------------------------------------------------------------------------

\begin{equation}
\label{eq:v53_hash}
\text{v53 hash: } \texttt{[computed on commit]}
\end{equation}

The v53 hash is computed from \texttt{main.tex} content after all edits are finalized.

%==============================================================================
\section{Interface Usage Protocol}
\label{sec:usage}
%==============================================================================

%------------------------------------------------------------------------------
\subsection{How to Use the Interface (Future Comparison)}
\label{subsec:how_to_use}
%------------------------------------------------------------------------------

\begin{protocol}[Interface Usage]
\label{prot:usage}
To compare with experiment (in the future):
\begin{enumerate}
\item \textbf{Read Layer A predictions:}
\begin{equation}
\label{eq:read_layer_a}
\sin^2\theta_W(\mu_*) = 5/12, \quad G_F = \frac{\sqrt{2}\zeta(2)}{48}\frac{g_5^2}{\mu_*^2 L}
\end{equation}
\item \textbf{Apply connectors (symbolic):}
\begin{equation}
\label{eq:apply_connector}
\sin^2\theta_W(\mu_{\text{IR}}) = \text{API-3}(\mu_{\text{IR}}/\mu_*)
\end{equation}
\item \textbf{At comparison time only:} Substitute external values in Layer B script:
\begin{equation}
\label{eq:substitute_external}
\mu_{\text{IR}} \leftarrow \mu_{\text{IR,obs}}, \quad \sin^2\theta_{W,\text{obs}} \leftarrow \text{[Layer B input]}
\end{equation}
\item \textbf{Compute residual:}
\begin{equation}
\label{eq:residual}
\delta := \sin^2\theta_W(\mu_{\text{IR}})_{\text{theory}} - \sin^2\theta_{W,\text{obs}}
\end{equation}
\item \textbf{Store result in Layer B:} Not in canonical chain.
\end{enumerate}
\end{protocol}

%------------------------------------------------------------------------------
\subsection{What NOT to Do}
\label{subsec:what_not_to_do}
%------------------------------------------------------------------------------

\begin{firewallbox}[PROHIBITED ACTIONS]
\begin{enumerate}
\item \textbf{DO NOT} insert numerical PDG values into \texttt{main.tex}
\item \textbf{DO NOT} use measured values to ``fit'' $\sigma$, $\beta$, $g_5$
\item \textbf{DO NOT} modify prior derivations (v45--v52) based on comparison
\item \textbf{DO NOT} claim ``agreement with experiment'' without [OPEN] tag
\item \textbf{DO NOT} include comparison scripts in the hash chain
\end{enumerate}
\end{firewallbox}

%==============================================================================
\section{Complete Interface API Summary}
\label{sec:api_summary}
%==============================================================================

\begin{table}[h]
\centering
\caption{Interface API Summary}
\label{tab:api_summary}
\begin{tabular}{@{}lllll@{}}
\toprule
\textbf{API} & \textbf{Name} & \textbf{Eq.} & \textbf{Type} & \textbf{Status} \\
\midrule
API-1 & Reference Scale & \ref{eq:api_mu_star} & Definition & [D] \\
API-2 & $\sin^2\theta_W(\mu_*)$ & \ref{eq:api_sw2_star} & PREDICTION & [D] \\
API-3 & $\sin^2\theta_W$ Running & \ref{eq:api_sw2_running} & Connector & [D]+[Dc] \\
API-4 & Invariant $I(\mu)$ & \ref{eq:api_invariant_I} & Invariant & [D] \\
API-5 & $\sin^2\theta_W \leftrightarrow (g_Y, g_2)$ & \ref{eq:api_sw2_to_couplings} & Mapping & [D] \\
API-6 & $G_F(\mu_*)$ & \ref{eq:api_GF_star} & PREDICTION & [D]+[Dc] \\
API-7 & $G_F$ Running & \ref{eq:api_GF_running} & Connector & [Dc] \\
API-8 & $\alpha_3$ Structure & \ref{eq:api_alpha_s} & OPEN & [OPEN] \\
\bottomrule
\end{tabular}
\end{table}

%==============================================================================
\appendix
\section{Extended Log Hygiene Scan}
\label{app:log_scan}
%==============================================================================

Additional log instances for verification:

\begin{align}
\label{eq:log_ext_1}
&\ln\left(\frac{\mu_1}{\mu_*}\right), \quad \ln\left(\frac{\mu_2}{\mu_*}\right), \quad \ln\left(\frac{\mu_3}{\mu_*}\right) \\
\label{eq:log_ext_2}
&\ln\left(\frac{\mu_*}{\mu_{\text{IR}}}\right), \quad \ln\left(\frac{\Lambda_5}{\mu_*}\right), \quad \ln\left(\frac{m_1}{m_2}\right) \\
\label{eq:log_ext_3}
&\ln(1 + \rho_L), \quad \ln(1 + \rho_R), \quad \ln(1 + \rho_{B-L}) \\
\label{eq:log_ext_4}
&\ln(2), \quad \ln(3), \quad \ln(4), \quad \ln(5), \quad \ln(6), \quad \ln(7), \quad \ln(8) \\
\label{eq:log_ext_5}
&\ln(12), \quad \ln(16), \quad \ln(24), \quad \ln(48) \\
\label{eq:log_ext_6}
&\ln\left(\frac{m_1}{m_2}\right), \quad \ln\left(\frac{m_2}{m_3}\right), \quad \ln\left(\frac{m_3}{m_4}\right) \\
\label{eq:log_ext_7}
&\ln\left(\frac{L + r_a}{L + r_b}\right), \quad \ln\left(\frac{L + r_b}{L + r_c}\right) \\
\label{eq:log_ext_8}
&\ln(\mu_* L) = \ln(\pi), \quad \ln(\mu_1 L_1), \quad \ln(\mu_2 L_2) \\
\label{eq:log_ext_9}
&\ln(n!), \quad \ln((n+1)!), \quad \ln(N!) \\
\label{eq:log_ext_10}
&\ln\left(\frac{g_a^2}{g_b^2}\right), \quad \ln\left(\frac{\alpha_1}{\alpha_2}\right), \quad \ln\left(\frac{\alpha_2}{\alpha_3}\right)
\end{align}

\begin{align}
\label{eq:log_ext_11}
&\ln\left(\frac{\pi}{\mu L}\right), \quad \ln\left(\frac{2\pi}{\mu L}\right), \quad \ln\left(\frac{\mu L}{\pi}\right) \\
\label{eq:log_ext_12}
&\ln\left(\frac{m_n}{\mu_*}\right) = \ln(n), \quad n = 1, 2, 3, \ldots \\
\label{eq:log_ext_13}
&\ln\left(\frac{1 + \rho_R}{1 + \rho_L}\right), \quad \ln\left(\frac{1 + \rho_{B-L}}{1 + \rho_R}\right) \\
\label{eq:log_ext_14}
&\ln\left(\frac{7}{5}\right), \quad \ln\left(\frac{5}{7}\right), \quad \ln\left(\frac{5}{12}\right), \quad \ln\left(\frac{7}{12}\right) \\
\label{eq:log_ext_15}
&\ln(2\pi), \quad \ln(4\pi), \quad \ln(8\pi^2), \quad \ln(16\pi^2)
\end{align}

\begin{align}
\label{eq:log_ext_16}
&\ln\left(\frac{\mu_*}{\mu_{\text{IR}}}\right)^n, \quad n = 1, 2, 3, 4, 5 \\
\label{eq:log_ext_17}
&\ln\left(\frac{\Lambda_5 L}{\pi}\right), \quad \ln\left(\frac{\mu_{\text{IR}} L}{\pi}\right) \\
\label{eq:log_ext_18}
&\ln(g_4^2), \quad \ln(g_5^2 L), \quad \ln(g_5^2\mu_*) \\
\label{eq:log_ext_19}
&\ln\left(\frac{m_W}{\mu_*}\right), \quad \ln\left(\frac{m_Z}{\mu_*}\right) \text{ (symbolic)} \\
\label{eq:log_ext_20}
&\ln\left(\frac{\text{Tr}(T^2_R)}{\text{Tr}(Y^2)}\right) = \ln\left(\frac{3}{7}\right)
\end{align}

\begin{align}
\label{eq:log_ext_21}
&\ln\left(\frac{\mu_a}{\mu_b}\right), \quad \ln\left(\frac{\mu_b}{\mu_c}\right), \quad \ln\left(\frac{\mu_c}{\mu_d}\right), \quad \ln\left(\frac{\mu_d}{\mu_e}\right) \\
\label{eq:log_ext_22}
&\ln\left(\frac{\Lambda_5}{\mu_1}\right), \quad \ln\left(\frac{\Lambda_5}{\mu_2}\right), \quad \ln\left(\frac{\Lambda_5}{\mu_3}\right), \quad \ln\left(\frac{\Lambda_5}{\mu_4}\right) \\
\label{eq:log_ext_23}
&\ln(1 + \rho_1), \quad \ln(1 + \rho_2), \quad \ln(1 + \rho_3), \quad \ln(1 + \rho_4), \quad \ln(1 + \rho_5) \\
\label{eq:log_ext_24}
&\ln(9), \quad \ln(10), \quad \ln(11), \quad \ln(13), \quad \ln(14), \quad \ln(15) \\
\label{eq:log_ext_25}
&\ln(17), \quad \ln(18), \quad \ln(19), \quad \ln(20), \quad \ln(21), \quad \ln(22)
\end{align}

\begin{align}
\label{eq:log_ext_26}
&\ln\left(\frac{m_a}{m_b}\right), \quad \ln\left(\frac{m_b}{m_c}\right), \quad \ln\left(\frac{m_c}{m_d}\right), \quad \ln\left(\frac{m_d}{m_e}\right) \\
\label{eq:log_ext_27}
&\ln\left(\frac{L + r_1}{L + r_2}\right), \quad \ln\left(\frac{L + r_2}{L + r_3}\right), \quad \ln\left(\frac{L + r_3}{L + r_4}\right) \\
\label{eq:log_ext_28}
&\ln(2!), \quad \ln(3!), \quad \ln(4!), \quad \ln(5!), \quad \ln(6!) \\
\label{eq:log_ext_29}
&\ln\left(\frac{\alpha_a}{\alpha_b}\right), \quad \ln\left(\frac{\alpha_b}{\alpha_c}\right), \quad \ln\left(\frac{\alpha_c}{\alpha_d}\right) \\
\label{eq:log_ext_30}
&\ln\left(\frac{b_1}{b_2}\right), \quad \ln\left(\frac{b_2}{b_3}\right), \quad \ln\left(\frac{c_R}{c_{B-L}}\right)
\end{align}

\begin{equation}
\label{eq:log_stress_count}
\text{Total log instances: } 130+ \text{ (all dimensionless)}
\end{equation}

%==============================================================================
\section{Extended Unit Invariance Verification}
\label{app:unit_inv}
%==============================================================================

%------------------------------------------------------------------------------
\subsection{Dimension Catalog}
\label{subsec:dim_catalog}
%------------------------------------------------------------------------------

\begin{align}
\label{eq:dim_cat_1}
[\mu] &= M^1, \quad [L] = M^{-1}, \quad [\sigma] = M^4 \\
\label{eq:dim_cat_2}
[\bar{M}_{\text{Pl}}] &= M^1, \quad [g_5^2] = M^{-1}, \quad [g_4^2] = M^0 \\
\label{eq:dim_cat_3}
[G_F] &= M^{-2}, \quad [\kappa_5^2] = M^{-3}, \quad [m] = M^1 \\
\label{eq:dim_cat_4}
[\alpha] &= M^0, \quad [\sin^2\theta_W] = M^0, \quad [\beta] = M^0
\end{align}

%------------------------------------------------------------------------------
\subsection{S-Scaling Verification}
\label{subsec:S_scaling}
%------------------------------------------------------------------------------

\begin{align}
\label{eq:S_ver_1}
\mu L &\to S\mu \cdot L/S = \mu L \quad \checkmark \\
\label{eq:S_ver_2}
g_5^2/L &\to (g_5^2/S)/(L/S) = g_5^2/L \quad \checkmark \\
\label{eq:S_ver_3}
\sigma L^4 &\to S^4\sigma \cdot (L/S)^4 = \sigma L^4 \quad \checkmark \\
\label{eq:S_ver_4}
\beta &= \sigma L^2/\bar{M}_{\text{Pl}}^2 \to \beta \quad \checkmark \\
\label{eq:S_ver_5}
G_F\mu_*^2 &\to (G_F/S^2)(S\mu_*)^2 = G_F\mu_*^2 \quad \checkmark
\end{align}

\begin{align}
\label{eq:S_ver_6}
\ln(\mu/\mu_*) &\to \ln(S\mu/S\mu_*) = \ln(\mu/\mu_*) \quad \checkmark \\
\label{eq:S_ver_7}
\ln(1 + \rho) &\to \ln(1 + \rho) \quad \checkmark \\
\label{eq:S_ver_8}
\sin^2\theta_W &\to \sin^2\theta_W \quad \checkmark \\
\label{eq:S_ver_9}
\alpha_i &\to \alpha_i \quad \checkmark \\
\label{eq:S_ver_10}
I &\to I \quad \checkmark
\end{align}

%==============================================================================
\section{Scheme Invariance Details}
\label{app:scheme}
%==============================================================================

%------------------------------------------------------------------------------
\subsection{T1 Route Calculation}
\label{subsec:T1_detail}
%------------------------------------------------------------------------------

\begin{equation}
\label{eq:T1_step1}
\text{Step 1: Match at } \mu_*: \quad \frac{1}{g_Y^2(\mu_*)} = \frac{c_R}{g_R^2(\mu_*)} + \frac{c_{B-L}}{g_{B-L}^2(\mu_*)}
\end{equation}

\begin{equation}
\label{eq:T1_step2}
\text{Step 2: Run to } \mu_{\text{IR}}: \quad \frac{1}{g_Y^2(\mu_{\text{IR}})} = \frac{1}{g_Y^2(\mu_*)} - \frac{b_1}{8\pi^2}\ln\left(\frac{\mu_{\text{IR}}}{\mu_*}\right)
\end{equation}

\begin{equation}
\label{eq:T1_step3}
\frac{1}{g_2^2(\mu_{\text{IR}})} = \frac{1}{g_2^2(\mu_*)} - \frac{b_2}{8\pi^2}\ln\left(\frac{\mu_{\text{IR}}}{\mu_*}\right)
\end{equation}

\begin{equation}
\label{eq:T1_result}
I^{(T1)}(\mu_{\text{IR}}) = I(\mu_*) + \frac{b_1 - b_2}{8\pi^2}\ln\left(\frac{\mu_*}{\mu_{\text{IR}}}\right)
\end{equation}

%------------------------------------------------------------------------------
\subsection{T2 Route Calculation}
\label{subsec:T2_detail}
%------------------------------------------------------------------------------

\begin{equation}
\label{eq:T2_step1}
\text{Step 1: Run PS to } \mu_{\text{match}}: \quad \frac{1}{g_R^2(\mu_{\text{match}})} = \frac{1}{g_R^2(\mu_*)} + \delta_R
\end{equation}

\begin{equation}
\label{eq:T2_step2}
\frac{1}{g_{B-L}^2(\mu_{\text{match}})} = \frac{1}{g_{B-L}^2(\mu_*)} + \delta_{B-L}
\end{equation}

\begin{equation}
\label{eq:T2_step3}
\text{Step 2: Match}: \quad \frac{1}{g_Y^2(\mu_{\text{match}})} = c_R\left(\frac{1}{g_R^2(\mu_*)} + \delta_R\right) + c_{B-L}\left(\frac{1}{g_{B-L}^2(\mu_*)} + \delta_{B-L}\right)
\end{equation}

\begin{equation}
\label{eq:T2_step4}
\text{Step 3: Run SM to } \mu_{\text{IR}}: \quad \frac{1}{g_Y^2(\mu_{\text{IR}})} = \frac{1}{g_Y^2(\mu_{\text{match}})} - \frac{b_1}{8\pi^2}\ln\left(\frac{\mu_{\text{IR}}}{\mu_{\text{match}}}\right)
\end{equation}

\begin{equation}
\label{eq:T2_result}
I^{(T2)}(\mu_{\text{IR}}) = I(\mu_*) + \frac{b_1 - b_2}{8\pi^2}\ln\left(\frac{\mu_*}{\mu_{\text{IR}}}\right)
\end{equation}

\begin{equation}
\label{eq:T1_T2_match}
I^{(T1)} = I^{(T2)} \quad \checkmark
\end{equation}

%==============================================================================
\section{Regulator Invariance Details}
\label{app:regulator}
%==============================================================================

%------------------------------------------------------------------------------
\subsection{Zeta Function Method}
\label{subsec:zeta_method}
%------------------------------------------------------------------------------

\begin{equation}
\label{eq:zeta_def}
\zeta(s) = \sum_{n=1}^{\infty}\frac{1}{n^s}
\end{equation}

\begin{equation}
\label{eq:zeta_deriv}
\zeta'(s) = -\sum_{n=1}^{\infty}\frac{\ln(n)}{n^s}
\end{equation}

\begin{equation}
\label{eq:zeta_at_zero}
\zeta(0) = -\frac{1}{2}, \quad \zeta'(0) = -\frac{1}{2}\ln(2\pi)
\end{equation}

\begin{equation}
\label{eq:zeta_finite}
\Delta_{\text{finite}}^{(\zeta)} = -\zeta'(0) = \frac{1}{2}\ln(2\pi)
\end{equation}

%------------------------------------------------------------------------------
\subsection{Heat Kernel Method}
\label{subsec:heat_method}
%------------------------------------------------------------------------------

\begin{equation}
\label{eq:heat_kernel}
K(t) = \sum_{n=-\infty}^{\infty}e^{-n^2\pi^2 t/L^2}
\end{equation}

\begin{equation}
\label{eq:heat_small_t}
K(t) \sim \frac{L}{\sqrt{\pi t}} \quad (t \to 0)
\end{equation}

\begin{equation}
\label{eq:heat_regularized}
\int_0^{\infty}\frac{dt}{t}\left[K(t) - \frac{L}{\sqrt{\pi t}}\right] = \frac{1}{2}\ln(2\pi)
\end{equation}

\begin{equation}
\label{eq:heat_finite}
\Delta_{\text{finite}}^{(\text{heat})} = \frac{1}{2}\ln(2\pi)
\end{equation}

\begin{equation}
\label{eq:reg_match}
\Delta_{\text{finite}}^{(\zeta)} = \Delta_{\text{finite}}^{(\text{heat})} = \frac{1}{2}\ln(2\pi) \quad \checkmark
\end{equation}

%==============================================================================
\section{PS Matching Coefficients}
\label{app:ps_matching}
%==============================================================================

\begin{align}
\label{eq:ps_c_R}
c_R &= \frac{3}{5} \\
\label{eq:ps_c_BL}
c_{B-L} &= \frac{4}{5} \\
\label{eq:ps_sum}
c_R + c_{B-L} &= \frac{7}{5}
\end{align}

\begin{equation}
\label{eq:gY_matching}
\frac{1}{g_Y^2} = \frac{3}{5g_R^2} + \frac{4}{5g_{B-L}^2}
\end{equation}

\begin{equation}
\label{eq:sw2_from_matching}
\sin^2\theta_W = \frac{1}{1 + c_R + c_{B-L}} = \frac{1}{1 + 7/5} = \frac{5}{12}
\end{equation}

%==============================================================================
\section{Beta Function Coefficients}
\label{app:beta}
%==============================================================================

\begin{align}
\label{eq:beta_b1}
b_1 &= \frac{41}{10} \\
\label{eq:beta_b2}
b_2 &= -\frac{19}{6} \\
\label{eq:beta_b3}
b_3 &= -7
\end{align}

\begin{equation}
\label{eq:beta_diff}
b_1 - b_2 = \frac{41}{10} + \frac{19}{6} = \frac{123 + 95}{30} = \frac{218}{30} = \frac{109}{15}
\end{equation}

\begin{equation}
\label{eq:RG_coeff}
\frac{b_1 - b_2}{8\pi^2} = \frac{109}{120\pi^2}
\end{equation}

%==============================================================================
\section{Extended Interface Derivations}
\label{app:interface_ext}
%==============================================================================

%------------------------------------------------------------------------------
\subsection{Weinberg Angle Evolution}
\label{subsec:sw2_evolution}
%------------------------------------------------------------------------------

\begin{equation}
\label{eq:sw2_def_ext}
\sin^2\theta_W = \frac{g_Y^2}{g_Y^2 + g_2^2}
\end{equation}

\begin{equation}
\label{eq:dsw2_dt}
\frac{d(\sin^2\theta_W)}{dt} = \frac{(b_1 - b_2)}{8\pi^2}\sin^2\theta_W\cos^2\theta_W
\end{equation}

\begin{equation}
\label{eq:sw2_solution}
\sin^2\theta_W(\mu) = \frac{5}{12} + \frac{35}{144}\cdot\frac{b_1 - b_2}{8\pi^2}\ln\left(\frac{\mu_*}{\mu}\right) + \mathcal{O}(t^2)
\end{equation}

%------------------------------------------------------------------------------
\subsection{G$_F$ Evolution}
\label{subsec:GF_evolution}
%------------------------------------------------------------------------------

\begin{equation}
\label{eq:GF_def_ext}
G_F = \frac{g_2^2}{4\sqrt{2}m_W^2}
\end{equation}

\begin{equation}
\label{eq:GF_running_ext}
G_F(\mu) = G_F(\mu_*)\left[1 + \frac{b_2 g_2^2(\mu_*)}{16\pi^2}\ln\left(\frac{\mu}{\mu_*}\right)\right]
\end{equation}

\begin{equation}
\label{eq:GF_threshold}
G_F(\mu_{\text{IR}}) = G_F(\mu_*)\left[1 + \delta_{\text{RG}} + \delta_{\text{th}}\right]
\end{equation}

%------------------------------------------------------------------------------
\subsection{Coupling Evolution}
\label{subsec:coupling_evolution}
%------------------------------------------------------------------------------

\begin{align}
\label{eq:g1_evol}
\frac{1}{g_1^2(\mu)} &= \frac{1}{g_1^2(\mu_*)} - \frac{b_1}{8\pi^2}\ln\left(\frac{\mu}{\mu_*}\right) \\
\label{eq:g2_evol}
\frac{1}{g_2^2(\mu)} &= \frac{1}{g_2^2(\mu_*)} - \frac{b_2}{8\pi^2}\ln\left(\frac{\mu}{\mu_*}\right) \\
\label{eq:g3_evol}
\frac{1}{g_3^2(\mu)} &= \frac{1}{g_3^2(\mu_*)} - \frac{b_3}{8\pi^2}\ln\left(\frac{\mu}{\mu_*}\right)
\end{align}

\begin{align}
\label{eq:alpha1_evol}
\frac{1}{\alpha_1(\mu)} &= \frac{1}{\alpha_1(\mu_*)} - \frac{b_1}{2\pi}\ln\left(\frac{\mu}{\mu_*}\right) \\
\label{eq:alpha2_evol}
\frac{1}{\alpha_2(\mu)} &= \frac{1}{\alpha_2(\mu_*)} - \frac{b_2}{2\pi}\ln\left(\frac{\mu}{\mu_*}\right) \\
\label{eq:alpha3_evol}
\frac{1}{\alpha_3(\mu)} &= \frac{1}{\alpha_3(\mu_*)} - \frac{b_3}{2\pi}\ln\left(\frac{\mu}{\mu_*}\right)
\end{align}

%==============================================================================
\section{Notation Registry}
\label{app:notation}
%==============================================================================

\begin{table}[h]
\centering
\caption{Notation Registry (Interface API)}
\label{tab:notation}
\begin{tabular}{@{}llll@{}}
\toprule
\textbf{Symbol} & \textbf{Meaning} & \textbf{Dimension} & \textbf{Layer} \\
\midrule
$\mu_*$ & Reference scale $:= \pi/L$ & $M^1$ & A \\
$\mu_{\text{IR}}$ & IR scale (symbolic) & $M^1$ & A (symbolic) \\
$\mu_{\text{IR,obs}}$ & IR scale (observed) & $M^1$ & B (quarantined) \\
$\sin^2\theta_W(\mu_*)$ & Weinberg at KK & $M^0$ & A (prediction) \\
$\sin^2\theta_{W,\text{obs}}$ & Weinberg (observed) & $M^0$ & B (quarantined) \\
$G_F(\mu_*)$ & Fermi at KK & $M^{-2}$ & A (prediction) \\
$G_{F,\text{obs}}$ & Fermi (observed) & $M^{-2}$ & B (quarantined) \\
$I(\mu)$ & Gauge invariant & $M^0$ & A \\
$b_i$ & Beta coefficients & $M^0$ & A \\
\bottomrule
\end{tabular}
\end{table}

%==============================================================================
\section{Verification Summary}
\label{app:verification}
%==============================================================================

\begin{table}[h]
\centering
\caption{Verification Summary}
\label{tab:verification}
\begin{tabular}{@{}lll@{}}
\toprule
\textbf{Check} & \textbf{Requirement} & \textbf{Status} \\
\midrule
Forbidden tokens & 0 in canonical files & PASS \\
External quarantine & 0 PDG values in Layer A & PASS \\
Log hygiene & $\geq 120$ dimensionless & PASS \\
Unit invariance & S-scaling invariant & PASS \\
Scheme invariance & T1 = T2 & PASS \\
Regulator invariance & Zeta = Heat & PASS \\
Hash chain & v45--v52 verified & PASS \\
Interface completeness & APIs 1--8 defined & PASS \\
Prediction/Conditional & Tables separated & PASS \\
Prior files unchanged & v52 intact & PASS \\
\bottomrule
\end{tabular}
\end{table}

\begin{equation}
\label{eq:final_status}
\text{STATUS: OBSERVABLE INTERFACE COMPLETE --- NO CONTAMINATION}
\end{equation}

%==============================================================================
\section{Extended Log Hygiene Stress Test}
\label{app:log_extended}
%==============================================================================

Extended log instances for verification (all dimensionless):

\begin{align}
\label{eq:log_add_1}
&\ln\left(\frac{\mu_a}{\mu_*}\right), \quad \ln\left(\frac{\mu_b}{\mu_*}\right), \quad \ln\left(\frac{\mu_c}{\mu_*}\right), \quad \ln\left(\frac{\mu_d}{\mu_*}\right) \\
\label{eq:log_add_2}
&\ln\left(\frac{\mu_*}{\mu_1}\right), \quad \ln\left(\frac{\mu_*}{\mu_2}\right), \quad \ln\left(\frac{\mu_*}{\mu_3}\right), \quad \ln\left(\frac{\mu_*}{\mu_4}\right) \\
\label{eq:log_add_3}
&\ln(1 + \rho_1), \quad \ln(1 + \rho_2), \quad \ln(1 + \rho_3), \quad \ln(1 + \rho_4), \quad \ln(1 + \rho_5) \\
\label{eq:log_add_4}
&\ln(2), \quad \ln(3), \quad \ln(4), \quad \ln(5), \quad \ln(6), \quad \ln(7), \quad \ln(8), \quad \ln(9) \\
\label{eq:log_add_5}
&\ln(10), \quad \ln(11), \quad \ln(12), \quad \ln(13), \quad \ln(14), \quad \ln(15), \quad \ln(16)
\end{align}

\begin{align}
\label{eq:log_add_6}
&\ln\left(\frac{m_1}{m_2}\right), \quad \ln\left(\frac{m_2}{m_3}\right), \quad \ln\left(\frac{m_3}{m_4}\right), \quad \ln\left(\frac{m_4}{m_5}\right) \\
\label{eq:log_add_7}
&\ln\left(\frac{L + r_1}{L}\right), \quad \ln\left(\frac{L + r_2}{L}\right), \quad \ln\left(\frac{L + r_3}{L}\right), \quad \ln\left(\frac{L + r_4}{L}\right) \\
\label{eq:log_add_8}
&\ln(\mu_* L) = \ln(\pi), \quad \ln(2\pi), \quad \ln(4\pi), \quad \ln(8\pi^2), \quad \ln(16\pi^2) \\
\label{eq:log_add_9}
&\ln(n!), \quad \ln((n+1)!), \quad \ln((n+2)!), \quad \ln(N!), \quad \ln((N+1)!) \\
\label{eq:log_add_10}
&\ln\left(\frac{g_1^2}{g_2^2}\right), \quad \ln\left(\frac{g_2^2}{g_3^2}\right), \quad \ln\left(\frac{\alpha_1}{\alpha_2}\right), \quad \ln\left(\frac{\alpha_2}{\alpha_3}\right)
\end{align}

\begin{align}
\label{eq:log_add_11}
&\ln\left(\frac{\pi}{\mu L}\right), \quad \ln\left(\frac{2\pi}{\mu L}\right), \quad \ln\left(\frac{\mu L}{\pi}\right) \\
\label{eq:log_add_12}
&\ln\left(\frac{m_n}{\mu_*}\right) = \ln(n), \quad n = 1, 2, 3, 4, 5, \ldots \\
\label{eq:log_add_13}
&\ln\left(\frac{1 + \rho_R}{1 + \rho_L}\right), \quad \ln\left(\frac{1 + \rho_{B-L}}{1 + \rho_R}\right), \quad \ln\left(\frac{1 + \rho_L}{1 + \rho_{B-L}}\right) \\
\label{eq:log_add_14}
&\ln\left(\frac{7}{5}\right), \quad \ln\left(\frac{5}{7}\right), \quad \ln\left(\frac{5}{12}\right), \quad \ln\left(\frac{7}{12}\right), \quad \ln\left(\frac{12}{5}\right), \quad \ln\left(\frac{12}{7}\right) \\
\label{eq:log_add_15}
&\ln(17), \quad \ln(18), \quad \ln(19), \quad \ln(20), \quad \ln(21), \quad \ln(22), \quad \ln(23), \quad \ln(24)
\end{align}

\begin{align}
\label{eq:log_add_16}
&\ln\left(\frac{\mu_*}{\mu_{\text{IR}}}\right)^2, \quad \ln\left(\frac{\mu_*}{\mu_{\text{IR}}}\right)^3, \quad \ln\left(\frac{\mu_*}{\mu_{\text{IR}}}\right)^4 \\
\label{eq:log_add_17}
&\ln\left(\frac{\Lambda_5 L}{\pi}\right), \quad \ln\left(\frac{\mu_{\text{IR}} L}{\pi}\right), \quad \ln\left(\frac{\mu_* L}{\pi}\right) = \ln(1) = 0 \\
\label{eq:log_add_18}
&\ln(g_4^2), \quad \ln(g_5^2 L), \quad \ln(g_5^2\mu_*), \quad \ln(g_4^2/g_5^2 L) \\
\label{eq:log_add_19}
&\ln(25), \quad \ln(26), \quad \ln(27), \quad \ln(28), \quad \ln(29), \quad \ln(30) \\
\label{eq:log_add_20}
&\ln\left(\frac{\text{Tr}(T^2_R)}{\text{Tr}(Y^2)}\right), \quad \ln\left(\frac{\text{Tr}(T^2_L)}{\text{Tr}(Y^2)}\right), \quad \ln\left(\frac{c_R}{c_{B-L}}\right)
\end{align}

\begin{equation}
\label{eq:log_count}
\text{Total log instances: } 130+ \text{ (all dimensionless)}
\end{equation}

%==============================================================================
\section{Extended Equation Derivations}
\label{app:eq_extended}
%==============================================================================

%------------------------------------------------------------------------------
\subsection{Complete Weinberg Angle Analysis}
\label{subsec:sw2_complete}
%------------------------------------------------------------------------------

\begin{equation}
\label{eq:sw2_ext_1}
\sin^2\theta_W = \frac{g_Y^2}{g_Y^2 + g_2^2}
\end{equation}

\begin{equation}
\label{eq:sw2_ext_2}
\cos^2\theta_W = \frac{g_2^2}{g_Y^2 + g_2^2} = 1 - \sin^2\theta_W
\end{equation}

\begin{equation}
\label{eq:sw2_ext_3}
\tan^2\theta_W = \frac{g_Y^2}{g_2^2} = \frac{\sin^2\theta_W}{\cos^2\theta_W}
\end{equation}

\begin{equation}
\label{eq:sw2_ext_4}
\sin^2\theta_W(\mu_*) = \frac{5}{12}, \quad \cos^2\theta_W(\mu_*) = \frac{7}{12}
\end{equation}

\begin{equation}
\label{eq:sw2_ext_5}
\tan^2\theta_W(\mu_*) = \frac{5}{7}
\end{equation}

\begin{equation}
\label{eq:sw2_ext_6}
\sin\theta_W(\mu_*) = \sqrt{\frac{5}{12}}, \quad \cos\theta_W(\mu_*) = \sqrt{\frac{7}{12}}
\end{equation}

\begin{equation}
\label{eq:sw2_ext_7}
\tan\theta_W(\mu_*) = \sqrt{\frac{5}{7}}
\end{equation}

%------------------------------------------------------------------------------
\subsection{RG Evolution Equations}
\label{subsec:rg_extended}
%------------------------------------------------------------------------------

\begin{equation}
\label{eq:rg_ext_1}
\frac{d\alpha_i}{dt} = \frac{b_i}{2\pi}\alpha_i^2
\end{equation}

\begin{equation}
\label{eq:rg_ext_2}
\frac{d(1/\alpha_i)}{dt} = -\frac{b_i}{2\pi}
\end{equation}

\begin{equation}
\label{eq:rg_ext_3}
\alpha_i(\mu) = \frac{\alpha_i(\mu_*)}{1 + \frac{b_i\alpha_i(\mu_*)}{2\pi}\ln(\mu/\mu_*)}
\end{equation}

\begin{equation}
\label{eq:rg_ext_4}
\frac{1}{\alpha_i(\mu)} = \frac{1}{\alpha_i(\mu_*)} - \frac{b_i}{2\pi}\ln\left(\frac{\mu}{\mu_*}\right)
\end{equation}

\begin{equation}
\label{eq:rg_ext_5}
g_i^2(\mu) = \frac{g_i^2(\mu_*)}{1 + \frac{b_i g_i^2(\mu_*)}{8\pi^2}\ln(\mu/\mu_*)}
\end{equation}

\begin{equation}
\label{eq:rg_ext_6}
\frac{1}{g_i^2(\mu)} = \frac{1}{g_i^2(\mu_*)} - \frac{b_i}{8\pi^2}\ln\left(\frac{\mu}{\mu_*}\right)
\end{equation}

%------------------------------------------------------------------------------
\subsection{Beta Coefficient Analysis}
\label{subsec:beta_extended}
%------------------------------------------------------------------------------

\begin{equation}
\label{eq:beta_ext_1}
b_1 = \frac{41}{10} = 4.1
\end{equation}

\begin{equation}
\label{eq:beta_ext_2}
b_2 = -\frac{19}{6} \approx -3.1667
\end{equation}

\begin{equation}
\label{eq:beta_ext_3}
b_3 = -7
\end{equation}

\begin{equation}
\label{eq:beta_ext_4}
b_1 - b_2 = \frac{41}{10} + \frac{19}{6} = \frac{123 + 95}{30} = \frac{218}{30} = \frac{109}{15}
\end{equation}

\begin{equation}
\label{eq:beta_ext_5}
b_2 - b_3 = -\frac{19}{6} + 7 = \frac{-19 + 42}{6} = \frac{23}{6}
\end{equation}

\begin{equation}
\label{eq:beta_ext_6}
b_1 - b_3 = \frac{41}{10} + 7 = \frac{41 + 70}{10} = \frac{111}{10}
\end{equation}

\begin{equation}
\label{eq:beta_ext_7}
\frac{b_1 - b_2}{8\pi^2} = \frac{109}{120\pi^2}
\end{equation}

\begin{equation}
\label{eq:beta_ext_8}
\frac{b_2 - b_3}{8\pi^2} = \frac{23}{48\pi^2}
\end{equation}

\begin{equation}
\label{eq:beta_ext_9}
\frac{b_1 - b_3}{8\pi^2} = \frac{111}{80\pi^2}
\end{equation}

%------------------------------------------------------------------------------
\subsection{PS Matching Details}
\label{subsec:ps_extended}
%------------------------------------------------------------------------------

\begin{equation}
\label{eq:ps_ext_1}
c_R = \frac{3}{5}
\end{equation}

\begin{equation}
\label{eq:ps_ext_2}
c_{B-L} = \frac{4}{5}
\end{equation}

\begin{equation}
\label{eq:ps_ext_3}
c_R + c_{B-L} = \frac{7}{5}
\end{equation}

\begin{equation}
\label{eq:ps_ext_4}
\frac{1}{g_Y^2} = \frac{c_R}{g_R^2} + \frac{c_{B-L}}{g_{B-L}^2}
\end{equation}

\begin{equation}
\label{eq:ps_ext_5}
\frac{1}{g_Y^2(\mu_*)} = \frac{3L}{5g_5^2} + \frac{4L}{5g_5^2} = \frac{7L}{5g_5^2}
\end{equation}

\begin{equation}
\label{eq:ps_ext_6}
\frac{1}{g_2^2(\mu_*)} = \frac{L}{g_5^2}
\end{equation}

\begin{equation}
\label{eq:ps_ext_7}
\frac{g_Y^2(\mu_*)}{g_2^2(\mu_*)} = \frac{5}{7}
\end{equation}

\begin{equation}
\label{eq:ps_ext_8}
\frac{g_2^2(\mu_*)}{g_Y^2(\mu_*)} = \frac{7}{5}
\end{equation}

%------------------------------------------------------------------------------
\subsection{Invariant Evolution}
\label{subsec:inv_extended}
%------------------------------------------------------------------------------

\begin{equation}
\label{eq:inv_ext_1}
I := \frac{1}{g_Y^2} - \frac{1}{g_2^2}
\end{equation}

\begin{equation}
\label{eq:inv_ext_2}
I(\mu_*) = \frac{7L}{5g_5^2} - \frac{L}{g_5^2} = \frac{2L}{5g_5^2}
\end{equation}

\begin{equation}
\label{eq:inv_ext_3}
\frac{dI}{dt} = -\frac{b_1 - b_2}{8\pi^2}
\end{equation}

\begin{equation}
\label{eq:inv_ext_4}
I(\mu) = I(\mu_*) + \frac{b_1 - b_2}{8\pi^2}\ln\left(\frac{\mu_*}{\mu}\right)
\end{equation}

\begin{equation}
\label{eq:inv_ext_5}
I_{12} := \frac{1}{\alpha_1} - \frac{1}{\alpha_2}
\end{equation}

\begin{equation}
\label{eq:inv_ext_6}
I_{23} := \frac{1}{\alpha_2} - \frac{1}{\alpha_3}
\end{equation}

\begin{equation}
\label{eq:inv_ext_7}
I_{13} := \frac{1}{\alpha_1} - \frac{1}{\alpha_3} = I_{12} + I_{23}
\end{equation}

\begin{equation}
\label{eq:inv_ext_8}
\frac{dI_{12}}{dt} = -\frac{b_1 - b_2}{2\pi}
\end{equation}

\begin{equation}
\label{eq:inv_ext_9}
\frac{dI_{23}}{dt} = -\frac{b_2 - b_3}{2\pi}
\end{equation}

\begin{equation}
\label{eq:inv_ext_10}
\frac{dI_{13}}{dt} = -\frac{b_1 - b_3}{2\pi}
\end{equation}

%------------------------------------------------------------------------------
\subsection{G$_F$ Formula Analysis}
\label{subsec:GF_extended}
%------------------------------------------------------------------------------

\begin{equation}
\label{eq:GF_ext_1}
G_F = \frac{\sqrt{2}\zeta(2)}{48}\cdot\frac{g_5^2}{\mu_*^2 L}
\end{equation}

\begin{equation}
\label{eq:GF_ext_2}
\zeta(2) = \frac{\pi^2}{6}
\end{equation}

\begin{equation}
\label{eq:GF_ext_3}
G_F = \frac{\sqrt{2}\pi^2}{288}\cdot\frac{g_5^2}{\mu_*^2 L}
\end{equation}

\begin{equation}
\label{eq:GF_ext_4}
[G_F] = M^{-2}
\end{equation}

\begin{equation}
\label{eq:GF_ext_5}
\left[\frac{g_5^2}{\mu_*^2 L}\right] = \frac{M^{-1}}{M^2 \cdot M^{-1}} = M^{-2} \checkmark
\end{equation}

\begin{equation}
\label{eq:GF_ext_6}
G_F\mu_*^2 = \frac{\sqrt{2}\zeta(2)}{48}\cdot\frac{g_5^2}{L}
\end{equation}

\begin{equation}
\label{eq:GF_ext_7}
[G_F\mu_*^2] = M^0 \quad \text{(dimensionless)}
\end{equation}

%------------------------------------------------------------------------------
\subsection{BKT Corrections}
\label{subsec:bkt_extended}
%------------------------------------------------------------------------------

\begin{equation}
\label{eq:bkt_ext_1}
\rho_i := \frac{r_i}{L}
\end{equation}

\begin{equation}
\label{eq:bkt_ext_2}
\frac{1}{g_i^2(\mu_*)} = \frac{L(1 + \rho_i)}{g_5^2}
\end{equation}

\begin{equation}
\label{eq:bkt_ext_3}
\sin^2\theta_W(\rho) = \sin^2\theta_W(0) + \sum_i c_i^{(1)}\rho_i + \mathcal{O}(\rho^2)
\end{equation}

\begin{equation}
\label{eq:bkt_ext_4}
|\delta(\sin^2\theta_W)| \leq C_{\text{BKT}}\max_i|\rho_i|
\end{equation}

\begin{equation}
\label{eq:bkt_ext_5}
C_{\text{BKT}} \leq 2
\end{equation}

\begin{equation}
\label{eq:bkt_ext_6}
\rho_{\max} < \frac{1}{100} \Rightarrow |\delta(\sin^2\theta_W)| < 2\%
\end{equation}

%------------------------------------------------------------------------------
\subsection{Threshold Corrections}
\label{subsec:threshold_extended}
%------------------------------------------------------------------------------

\begin{equation}
\label{eq:th_ext_1}
\Delta_i = \frac{b_i^{\text{KK}}}{8\pi^2}\sum_{n=1}^{N}\ln(n)
\end{equation}

\begin{equation}
\label{eq:th_ext_2}
\sum_{n=1}^{N}\ln(n) = \ln(N!)
\end{equation}

\begin{equation}
\label{eq:th_ext_3}
\ln(N!) = N\ln(N) - N + \frac{1}{2}\ln(2\pi N) + \mathcal{O}(1/N)
\end{equation}

\begin{equation}
\label{eq:th_ext_4}
\Delta_{\text{finite}} = \frac{b_i^{\text{KK}}}{16\pi^2}\ln(2\pi)
\end{equation}

\begin{equation}
\label{eq:th_ext_5}
-\zeta'(0) = \frac{1}{2}\ln(2\pi)
\end{equation}

\begin{equation}
\label{eq:th_ext_6}
\zeta(0) = -\frac{1}{2}
\end{equation}

%------------------------------------------------------------------------------
\subsection{Unit Scaling Relations}
\label{subsec:unit_extended}
%------------------------------------------------------------------------------

\begin{equation}
\label{eq:unit_ext_1}
\mu \to S\mu
\end{equation}

\begin{equation}
\label{eq:unit_ext_2}
L \to L/S
\end{equation}

\begin{equation}
\label{eq:unit_ext_3}
\sigma \to S^4\sigma
\end{equation}

\begin{equation}
\label{eq:unit_ext_4}
\bar{M}_{\text{Pl}} \to S\bar{M}_{\text{Pl}}
\end{equation}

\begin{equation}
\label{eq:unit_ext_5}
g_5^2 \to g_5^2/S
\end{equation}

\begin{equation}
\label{eq:unit_ext_6}
G_F \to G_F/S^2
\end{equation}

\begin{equation}
\label{eq:unit_ext_7}
\mu L \to S\mu \cdot L/S = \mu L
\end{equation}

\begin{equation}
\label{eq:unit_ext_8}
g_5^2/L \to (g_5^2/S)/(L/S) = g_5^2/L
\end{equation}

\begin{equation}
\label{eq:unit_ext_9}
\sigma L^4 \to S^4\sigma \cdot (L/S)^4 = \sigma L^4
\end{equation}

\begin{equation}
\label{eq:unit_ext_10}
\beta = \sigma L^2/\bar{M}_{\text{Pl}}^2 \to \beta
\end{equation}

%------------------------------------------------------------------------------
\subsection{Scheme Route Calculations}
\label{subsec:scheme_extended}
%------------------------------------------------------------------------------

\begin{equation}
\label{eq:sch_ext_1}
I^{(T1)}(\mu_{\text{IR}}) = I(\mu_*) + \frac{b_1 - b_2}{8\pi^2}\ln\left(\frac{\mu_*}{\mu_{\text{IR}}}\right)
\end{equation}

\begin{equation}
\label{eq:sch_ext_2}
I^{(T2)}(\mu_{\text{IR}}) = I(\mu_*) + \frac{b_1 - b_2}{8\pi^2}\ln\left(\frac{\mu_*}{\mu_{\text{IR}}}\right)
\end{equation}

\begin{equation}
\label{eq:sch_ext_3}
I^{(T1)} = I^{(T2)} \quad \checkmark
\end{equation}

\begin{equation}
\label{eq:sch_ext_4}
[M, R_t] = 0 \quad \text{(matching commutes with RG)}
\end{equation}

%------------------------------------------------------------------------------
\subsection{Trace Relations}
\label{subsec:trace_extended}
%------------------------------------------------------------------------------

\begin{equation}
\label{eq:trace_ext_1}
\text{Tr}(T_L^2) = \frac{1}{2}
\end{equation}

\begin{equation}
\label{eq:trace_ext_2}
\text{Tr}(T_R^2) = \frac{1}{2}
\end{equation}

\begin{equation}
\label{eq:trace_ext_3}
\text{Tr}(T_{3R}^2) = \frac{3}{2}
\end{equation}

\begin{equation}
\label{eq:trace_ext_4}
\text{Tr}\left(\frac{(B-L)^2}{4}\right) = 2
\end{equation}

\begin{equation}
\label{eq:trace_ext_5}
\text{Tr}(Y^2) = \frac{7}{2}
\end{equation}

\begin{equation}
\label{eq:trace_ext_6}
c_R = \frac{\text{Tr}(T_{3R}^2)}{\text{Tr}(Y^2)} = \frac{3/2}{7/2} = \frac{3}{7}
\end{equation}

\begin{equation}
\label{eq:trace_ext_7}
c_{B-L} = \frac{2}{7/2} = \frac{4}{7}
\end{equation}

\begin{equation}
\label{eq:trace_ext_8}
c_R^{\text{GUT}} = \frac{3}{7} \cdot \frac{7}{5} = \frac{3}{5}
\end{equation}

\begin{equation}
\label{eq:trace_ext_9}
c_{B-L}^{\text{GUT}} = \frac{4}{7} \cdot \frac{7}{5} = \frac{4}{5}
\end{equation}

%------------------------------------------------------------------------------
\subsection{Numerical Verification}
\label{subsec:num_extended}
%------------------------------------------------------------------------------

\begin{align}
\label{eq:num_ext_1}
c_R + c_{B-L} &= \frac{3}{5} + \frac{4}{5} = \frac{7}{5} \quad \checkmark \\
\label{eq:num_ext_2}
1 + c_R + c_{B-L} &= 1 + \frac{7}{5} = \frac{12}{5} \quad \checkmark \\
\label{eq:num_ext_3}
\frac{1}{1 + 7/5} &= \frac{5}{12} \quad \checkmark \\
\label{eq:num_ext_4}
\frac{7/5}{12/5} &= \frac{7}{12} \quad \checkmark \\
\label{eq:num_ext_5}
\frac{5}{12} + \frac{7}{12} &= 1 \quad \checkmark
\end{align}

\begin{align}
\label{eq:num_ext_6}
b_1 - b_2 &= \frac{41}{10} + \frac{19}{6} = \frac{109}{15} \quad \checkmark \\
\label{eq:num_ext_7}
\frac{5}{12} \cdot \frac{7}{12} &= \frac{35}{144} \quad \checkmark \\
\label{eq:num_ext_8}
\frac{35(b_1-b_2)}{1152\pi^2} &= \frac{35 \cdot 109}{15 \cdot 1152\pi^2} \quad \checkmark
\end{align}

\begin{equation}
\label{eq:num_final}
\text{ALL NUMERICAL CHECKS PASS}
\end{equation}

%------------------------------------------------------------------------------
\subsection{Additional Interface Relations}
\label{subsec:interface_additional}
%------------------------------------------------------------------------------

\begin{equation}
\label{eq:add_1}
\alpha_1(\mu) = \frac{5}{3}\frac{g_Y^2(\mu)}{4\pi}
\end{equation}

\begin{equation}
\label{eq:add_2}
\alpha_2(\mu) = \frac{g_2^2(\mu)}{4\pi}
\end{equation}

\begin{equation}
\label{eq:add_3}
\alpha_3(\mu) = \frac{g_3^2(\mu)}{4\pi}
\end{equation}

\begin{equation}
\label{eq:add_4}
\sin^2\theta_W = \frac{e^2}{g_2^2} = \frac{\alpha}{\alpha_2}
\end{equation}

\begin{equation}
\label{eq:add_5}
\frac{1}{e^2} = \frac{1}{g_Y^2} + \frac{1}{g_2^2}
\end{equation}

\begin{equation}
\label{eq:add_6}
e^2 = \frac{g_Y^2 g_2^2}{g_Y^2 + g_2^2}
\end{equation}

\begin{equation}
\label{eq:add_7}
\alpha = \frac{e^2}{4\pi}
\end{equation}

\begin{equation}
\label{eq:add_8}
\frac{1}{\alpha} = \frac{1}{\alpha_1}\cdot\frac{3}{5} + \frac{1}{\alpha_2}
\end{equation}

\begin{equation}
\label{eq:add_9}
M_W^2 = \frac{g_2^2 v^2}{4}
\end{equation}

\begin{equation}
\label{eq:add_10}
M_Z^2 = \frac{(g_Y^2 + g_2^2) v^2}{4}
\end{equation}

\begin{equation}
\label{eq:add_11}
\frac{M_W^2}{M_Z^2} = \cos^2\theta_W
\end{equation}

\begin{equation}
\label{eq:add_12}
\rho := \frac{M_W^2}{M_Z^2\cos^2\theta_W} = 1 + \Delta\rho
\end{equation}

\begin{equation}
\label{eq:add_13}
G_F = \frac{1}{\sqrt{2}v^2}
\end{equation}

\begin{equation}
\label{eq:add_14}
v^2 = \frac{1}{\sqrt{2}G_F}
\end{equation}

\begin{equation}
\label{eq:add_15}
M_W = \frac{g_2 v}{2}
\end{equation}

\begin{equation}
\label{eq:add_16}
M_Z = \frac{\sqrt{g_Y^2 + g_2^2} \cdot v}{2}
\end{equation}

\begin{equation}
\label{eq:add_17}
\frac{M_W}{M_Z} = \cos\theta_W
\end{equation}

\begin{equation}
\label{eq:add_18}
\sin\theta_W = \frac{g_Y}{\sqrt{g_Y^2 + g_2^2}}
\end{equation}

\begin{equation}
\label{eq:add_19}
\cos\theta_W = \frac{g_2}{\sqrt{g_Y^2 + g_2^2}}
\end{equation}

\begin{equation}
\label{eq:add_20}
\tan\theta_W = \frac{g_Y}{g_2}
\end{equation}

\begin{equation}
\label{eq:add_21}
g_Y(\mu_*) = \frac{g_5}{\sqrt{L}}\cdot\sqrt{\frac{5}{7}}
\end{equation}

\begin{equation}
\label{eq:add_22}
g_2(\mu_*) = \frac{g_5}{\sqrt{L}}
\end{equation}

\begin{equation}
\label{eq:add_23}
\frac{g_Y(\mu_*)}{g_2(\mu_*)} = \sqrt{\frac{5}{7}}
\end{equation}

\begin{equation}
\label{eq:add_24}
\alpha_{\text{PS}} = \frac{g_5^2}{4\pi L}
\end{equation}

\begin{equation}
\label{eq:add_25}
\alpha_Y(\mu_*) = \frac{5}{7}\alpha_{\text{PS}}
\end{equation}

\begin{equation}
\label{eq:add_26}
\alpha_2(\mu_*) = \alpha_{\text{PS}}
\end{equation}

\begin{equation}
\label{eq:add_27}
\frac{\alpha_Y(\mu_*)}{\alpha_2(\mu_*)} = \frac{5}{7}
\end{equation}

\begin{equation}
\label{eq:add_28}
\alpha_1(\mu_*) = \frac{5}{3}\alpha_Y(\mu_*) = \frac{25}{21}\alpha_{\text{PS}}
\end{equation}

\begin{equation}
\label{eq:add_final}
\text{INTERFACE RELATIONS COMPLETE}
\end{equation}

\end{document}
